\subsection{Set}
\textbf{Definition}\\
\texttt{set<type> s;}\\
\texttt{multiset<type> s;}\\
\texttt{unordered\_set<type> s;}\\
\texttt{unordered\_multiset<type, hasher> s;}\\
\texttt{hasher}: Refer Custom Hash Function.\\


\textbf{Set operations}\\
\begin{tabularx}{\linewidth}{lX}
\texttt{s.insert(item);} & \\
\texttt{s.erase(item);} & \\
\texttt{s.count(item);} & \\
\texttt{s.find(item);} & Returns the iterator that points to \texttt{item}.\\
\texttt{s.insert()} & Inserts elements in s.\\
& Eg: Insert elements from t to s.\\
& \texttt{v.insert(v.end(), u.begin(), u.end());}\\
\texttt{s.size()} & \\
& Not availabe for \texttt{unordered\_set}  and \texttt{unordered\_multiset}.\\
\end{tabularx}

\begin{tabularx}{\linewidth}{lX}
\texttt{s.lower\_bound(x)} & Returns iterator to the smallest element that is larger than or equal to x.\\
\texttt{s.upper\_bound(x)} & Returns iterator to the smallest element that is larger than x.\\
\end{tabularx}

\textcolor{red}{NOTE: erase method remove all instances of the item in the multiset.}

\subsection{Map}

\texttt{map} : Balanced binary tree.\\
\texttt{unordered\_map}: Hash list.\\

\textbf{Definition:}\\
\texttt{map<type1, type2> m;}\\
Eg: \texttt{map<string, int> m;}\\
\texttt{unordered\_map<type1, type2, hasher> m;}\\
\texttt{hasher}: Refer Custom Hash Function.\\

\textbf{Operations on map}\\
\begin{tabularx}{\linewidth}{lX}
\texttt{m["banana"]} & Retrive value associated with \texttt{banana}.\\
\texttt{m.count("banana")} & Return 1 if the key is present else 0.\\
\texttt{m.size();} & \\
\end{tabularx}
